% !TEX program = lualatex
% !TeX encoding = UTF-8
\PassOptionsToPackage{unicode}{hyperref}
\PassOptionsToPackage{bookmarksnumbered,bookmarksopen}{hyperref}
\documentclass[11pt,a4paper]{article}

% ============================================================
% إعدادات اللغة العربية
% ============================================================
\usepackage{polyglossia}
\setmainlanguage{arabic}
\setotherlanguage{english}

% خطوط عربية
\newfontfamily\arabicfont{Amiri}[Script=Arabic, Language=Arabic]
\newfontfamily\arabicfontsf{Amiri}[Script=Arabic, Language=Arabic]
\newfontfamily\arabicfonttt{Amiri}[Script=Arabic, Language=Arabic]
\newfontfamily\englishfont{Times New Roman}
\setmonofont{Latin Modern Mono}

% إزالة تحذيرات hyperref
\makeatletter
\let\@ensure@LTR\relax
\makeatother

% حزم الرياضيات والتنسيق
\usepackage{amsmath,amssymb,amsthm,mathtools,bm}
\usepackage{geometry}
\usepackage{enumitem}
\usepackage{siunitx}
\usepackage{booktabs}
\usepackage{array}
\usepackage{slashed}
\usepackage{graphicx}
\usepackage{caption}
\usepackage{subcaption}
\usepackage{braket}
\usepackage{tensor}
\usepackage{makecell}
\usepackage[most]{tcolorbox}
\usepackage{pgfplots}
\usepackage{float}
\usepackage{tikz}
\usepackage{multicol}
\usepackage{lipsum}
\usepackage{microtype}
\usepackage{hyperref}
\pgfplotsset{compat=1.18}
\usetikzlibrary{arrows.meta,positioning,shapes.geometric,fit,calc,
	decorations.pathreplacing,patterns,quotes,angles,
	shadows,decorations.markings}

\geometry{margin=1in}
\setlength{\emergencystretch}{3em}
\sloppy

% أوامر مختصرة (مطابقة للأصل)
\newcommand{\Keff}{K_{\mathrm{eff}}}
\newcommand{\kappac}{\kappa_c}
\newcommand{\be}{\beta}
\newcommand{\ga}{\gamma}
\newcommand{\lun}{\mathcal{L}}

% بيئات النظريات
\newtheorem{theorem}{مبرهنة}[section]
\newtheorem{lemma}{تمهيدية}[section]
\newtheorem{axiom}{مسلَمة}[section]
\newtheorem{remark}{ملاحظة}[section]
\newtheorem{proposition}{اقتراح}[section]
\newtheorem{definition}{تعريف}[section]
\newtheorem{assumption}{افتراض}[section]
\newtheorem{corollary}{نتيجة طبيعية}[section]

% بيانات PDF
\hypersetup{
	pdftitle={الملحق Climat01: علم المناخ التنسيقي لنظرية YUCT — القياس الشامل للتقلبات المناخية، السعة الحرارية للغلاف الجوي، التنسيق المحيطي والتنبؤ بالتحولات الحرجة},
	pdfauthor={أليكسي في. ياكوشيف},
	pdfsubject={كفاءة التنسيق (K_eff)، تقلبية المناخ، النشاط الشمسي، محتوى الحرارة المحيطي، توازن كتلة الأنهار الجليدية، التأثير البركاني، التنبؤات العمياء، التحولات الحرجة},
	pdfkeywords={YUCT, كفاءة التنسيق, تقلبية المناخ, النشاط الشمسي, محتوى الحرارة المحيطي, توازن كتلة الأنهار الجليدية, التأثير البركاني, التنبؤات العمياء, التحولات الحرجة, β = 0.67},
	breaklinks=true,
	pdfencoding=auto,
	bookmarksnumbered=true,
	bookmarksopen=true,
	bookmarksopenlevel=1
}

\title{الملحق Climat01: علم المناخ التنسيقي لنظرية YUCT\\
	\large القياس الشامل للتقلبات المناخية، السعة الحرارية للغلاف الجوي، التنسيق المحيطي والتنبؤ بالتحولات الحرجة}
\author{أليكسي في. ياكوشيف}
\date{مارس 2026}

\begin{document}
	\sloppy
	
	% صفحة العنوان
	\begin{titlepage}
		\centering
		\vspace*{1cm}
		{\Huge\textbf{الملحق Climat01: علم المناخ التنسيقي لنظرية YUCT}}\\[0.5cm]
		{\large\textit{القياس الشامل للتقلبات المناخية، السعة الحرارية للغلاف الجوي، التنسيق المحيطي والتنبؤ بالتحولات الحرجة}}\\[2cm]
		
		{\Large أليكسي في. ياكوشيف}\\[0.3cm]
		\large أبحاث ياكوشيف\\
		نظرية YUCT: \url{https://yuct.org/}\\
		بروتوكول YPSDC: \url{https://ypsdc.com/}\\[1cm]
		
		\textbf{DOI: \url{https://doi.org/10.5281/zenodo.18444598}}\\[0.5cm]
		
		نُشِرت نسخة مختصرة من هذا العمل سابقاً وهي متاحة على\\
		\url{https://doi.org/10.5281/zenodo.18362308}\\[0.5cm]
		مارس 2026 \quad V.2.1\\[2cm]
		
		\begin{minipage}{0.9\textwidth}
			\begin{abstract}
				\noindent
				يقدم هذا الملحق تطبيقاً شاملاً لنظرية التنسيق الموحّد لياكوشيف (YUCT) على النظام المناخي للأرض. نُبرهن أن قانون الخطأ القياسي الشامل \(\varepsilon = \kappa_c \alpha K_{\mathrm{eff}}^{-\beta}\) حيث \(\beta=0.67\) يحكم التقلبية المتعددة العقود لدرجة الحرارة العالمية، ومحتوى الحرارة المحيطية، وفقدان كتلة الصفائح الجليدية. يعمل النشاط الشمسي (أعداد وولف) كمؤشر خارجي يُعدّل كفاءة التنسيق \(K_{\mathrm{eff}}\) للنظام المناخي. باستخدام نوافذ متحركة مدتها 30 عاماً، نجد ميلاً ثابتاً لانحدار \(\ln \sigma_T\) على \(\ln W\) قدره \(-0.70\pm0.11\) (\(R^2=0.62\))، وهو في توافق ممتاز مع الأس النظري. يظهر نفس القياس في درجة حرارة سطح البحر (SST) ومحتوى الحرارة المحيطي، مع اختلافات إقليمية تُفسَّر بالقصور الذاتي المحيطي ودورة العقدة القمرية. بالنسبة للغلاف الجليدي، يكون أس القياس \(\approx 0.5\)، مما يعكس عمليات ذوبان غير خطية. قُدِّم نموذج جديد للتأثير البركاني، يربط احتمال الثوران بالنشاط المغناطيسي الأرضي عبر نفس الأس \(\beta=0.67\)، ويُدمج العمق البصري للهباء الجوي كحد توهين أُسّي في معادلة التقلب. باستخدام النموذج المتكامل، نُنتج تنبؤات عمياء لشذوذات درجة حرارة السطح لعام 2050، مع توزيعات إقليمية، ونناقش اقتراب النظام المناخي من عتبة حرجة (\(K_{\mathrm{eff}} \to K_{\mathrm{crit}}\)) حوالي عام 2045. يُوحِّد الإطار ديناميكيات المناخ مع الظواهر الاقتصادية والجيوفيزيائية تحت مبدأ تنسيقي واحد، مما يوفر تنبؤات قابلة للاختبار للعقود القادمة.
			\end{abstract}
		\end{minipage}
		
		\vfill
		\small
		\textbf{كلمات مفتاحية:} YUCT، كفاءة التنسيق، تقلبية المناخ، النشاط الشمسي، محتوى الحرارة المحيطي، توازن كتلة الأنهار الجليدية، التأثير البركاني، التنبؤات العمياء، التحولات الحرجة.
	\end{titlepage}
	
	\tableofcontents
	
	\section*{المسلَمة 0: الأولوية الوجودية للتنسيق}
	\addcontentsline{toc}{section}{المسلَمة 0: الأولوية الوجودية للتنسيق}
	
	كل نظرية تصف السكون. لكن التنسيق هو العملية التي تجعل السكون ممكناً.
	أي عبارة، أي نموذج، أي قانون لا يوجد إلا كعنصر مُفعَّل من قاموس.
	القاموس والفهرس لا يمكن اشتقاقهما من السكون الذي يصفانه — بل هما سابقان وجودياً.
	
	بروتوكول YPSDC ليس "إحدى النظريات".
	إنه البروتوكول الذي يكمن وراء أي نظرية، بما في ذلك تلك التي تصوغه.
	رفضه مستحيل دون رفض القدرة على الصياغة.
	
	بهذا المعنى، التنسيق ليس "أقوى" من النظريات الأخرى — إنه \textbf{أكثر أولوية}.
	السكون لا يُنتج؛ السكون مُنتَج.
	
	تُعتمد هذه المسلَمة كأساس وجودي جوهري لنظرية YUCT.
	
	\section{مقدمة}
	
	النظام المناخي للأرض هو نظام معقد عالي التنسيق. يتبادل الغلاف الجوي والمحيطات والغلاف الجليدي والمحيط الحيوي الطاقةَ والزخم والمادة باستمرار، ويتطلب الحفاظ على ظروف الهولوسين المستقرة درجة عالية من التنسيق. في إطار نظرية التنسيق الموحّد لياكوشيف (YUCT)، يُوصف أي نظام منسق بالثالوث \(D + I \cdot R\):
	
	\begin{itemize}
		\item \(D\) (القاموس) — القوانين الفيزيائية للهيدروديناميكا، النقل الإشعاعي، الديناميكا الحرارية؛ تُوزَّع مسبقاً.
		\item \(I\) (المعلومة) — متغيرات الحالة الراهنة: درجة الحرارة، الضغط، الملوحة، الغطاء الجليدي.
		\item \(R\) (الرنين) — أنماط التقلب المتسقة (ENSO, NAO, AMO) التي تزامن مناطق شاسعة.
	\end{itemize}
	
	تحدد كفاءة التنسيق \(\Keff\) للنظام المناخي سرعة عودته إلى التوازن بعد الاضطرابات. وفقاً لقانون YUCT الشامل، يتناسب الخطأ النسبي (سعة التقلبات) مع
	\[
	\epsilon = \kappac \alpha \Keff^{-\be}, \qquad \be = 0.67 \pm 0.03,\ \kappac = 0.33,
	\]
	حيث يمكن تفسير \(\epsilon\) على أنها السعة المعيارية للتقلبية المناخية (مثلاً، الانحراف المعياري لشذوذات درجة الحرارة) و\(\alpha\) ثابت خاص بالنظام.
	
	نُظهر في هذا الملحق أن هذا القانون ينطبق على درجة الحرارة العالمية، حيث يعمل النشاط الشمسي (أعداد وولف \(W\)) كمؤشر تنسيق خارجي. الأس الناتج يُطابق \(\be = 0.67\) ويتزامن مع الأس الذي وُجد سابقاً في الاقتصاد (الملحق F). وهذا يقدم تأكيداً قوياً متعدد التخصصات لعالمية YUCT. كما نوسع التحليل ليشمل المكون المحيطي والغلاف الجليدي والبنية القطاعية لدرجة حرارة الغلاف الجوي، مُنشئين روابط مع المد والجزر القمري والاعتماد الثقالي (الملحق P).
	
	\section{الإطار النظري}
	
	\subsection{كفاءة التنسيق للنظام المناخي}
	
	يعمل الغلاف الجوي كمحرك حراري يحول الطاقة الحرارية إلى طاقة حركية. يمكن التعبير عن كفاءته كالتالي
	\[
	\eta = \frac{T_h - T_c}{T_h},
	\]
	حيث \(T_h\) درجة حرارة المدخل الحراري (السطح المداري) و\(T_c\) درجة حرارة المخرج الحراري (التروبوسفير العلوي القطبي). في YUCT ترتبط كفاءة التنسيق الجوي بـ \(\eta\):
	\[
	\Keff^{\text{atm}} = K_0 \left( \frac{\eta}{\eta_0} \right)^{\be}.
	\]
	
	قدرة الغلاف الجوي على تخزين الطاقة الحرارية (السعة الحرارية الفعالة) تتغير مع \(\Keff\) وفقاً لـ
	\[
	C_{\text{heat}} = C_0 \cdot \Keff^{\ga}, \qquad \ga \approx 0.3 \pm 0.05,
	\]
	حيث \(\ga\) هو نفس الأس المستخدم في الملحق P، رابطاً بين درجة الحرارة وثابت الجاذبية.
	
	\subsection{قياس تقلبية درجة الحرارة}
	
	من القانون الشامل، سعة التقلب النسبي تتناسب عكسياً مع \(\Keff^{\be}\). بالنسبة لدرجة الحرارة العالمية:
	\[
	\frac{\sigma_T}{\langle T \rangle} \propto \Keff^{-\be}.
	\]
	بافتراض أن \(\Keff\) تتناسب طردياً مع النشاط الشمسي \(\langle W \rangle\) (مُوسَّطاً على زمن مميز)، نحصل على العلاقة القابلة للاختبار:
	\[
	\sigma_T \propto \langle W \rangle^{-\be}.
	\]
	
	\subsection{التأثير البشري في نموذج التنسيق}
	
	ينجم تغير المناخ الحديث عن عوامل طبيعية (النشاط الشمسي، البراكين) وبشرية (غازات الدفيئة، الهباء الجوي). في YUCT، يؤثر التأثير الخارجي على كفاءة التنسيق. ولمراعاة التأثير البشري، يُضاف حد يتناسب مع التأثير الإشعاعي \(\Delta F_{\mathrm{anth}}(t)\) إلى معادلة تطور \(K_{\mathrm{eff}}\):
	
	\[
	\frac{\partial K}{\partial t} = r K \left(1 - \frac{K}{K_{\mathrm{opt}}}\right) - \mu \kappa_c \alpha K^{1-\beta} 
	+ D\nabla^2 K + \gamma_{\mathrm{anth}} \frac{\Delta F_{\mathrm{anth}}(t)}{\Delta F_0} K^{1-\beta} + \xi(t),
	\]
	
	حيث \(\gamma_{\mathrm{anth}}\) حساسية كفاءة التنسيق للتأثير البشري (مُعايَرة من البيانات التاريخية). نستخدم في حساباتنا \(\Delta F_{\mathrm{anth}}(t)\) من سيناريوهات SSP2-4.5 و SSP5-8.5 (IPCC AR6). مساهمة التأثير البشري في انخفاض \(K_{\mathrm{eff}}\) قابلة للمقارنة مع تأثير النشاط الشمسي، مما يسمح بنمذجة كل من التذبذبات المتعددة العقود والاتجاه طويل الأمد.
	
	\subsection{النشاط البركاني كنتيجة لعمليات التنسيق}
	
	تُعدّل الثورات البركانية الكبيرة المناخ عبر الهباء الجوي الستراتوسفيري. لكن إحصائياتها ليست عشوائية بحتة: تُظهر السجلات التاريخية والمناخية القديمة ارتباطاً بين تواتر الثورات القوية ومراحل النشاط الشمسي والاضطرابات المغناطيسية الأرضية (مثلاً، دورة 11 عاماً وتوافقياتها). يُفسَّر هذا في YUCT بأن التغيرات في كفاءة التنسيق \(\Keff\) لا تؤثر فقط على الغلاف الجوي والمحيط بل أيضاً على العمليات الجيوديناميكية.
	
	\subsubsection{المؤشر المغناطيسي الأرضي كمُنبِّه للثوران}
	
	المؤشر المتكامل الذي يربط النشاط الشمسي بعمليات الأرض الصلبة هو \textbf{المؤشر المغناطيسي الأرضي \(Kp\)} (أو مشتقاته، مثلاً \(Ap\)). ترتبط فترات الاضطراب المغناطيسي الأرضي المتزايد (ذرى النشاط الشمسي) بعدد أكبر من الثورات الكبيرة، خصوصاً في المناطق الحساسة للإجهادات المدية.
	
	من قانون الخطأ القياسي الشامل، السعة النسبية للتغيرات المغناطيسية الأرضية \(\delta Kp / \langle Kp \rangle\) تتغير مع \(\Keff^{-\beta}\). يمكن حينئذ كتابة احتمال الثوران لكل وحدة زمن:
	
	\[
	P_{\text{eruption}}(t) = P_0 \cdot \left( \frac{Kp(t)}{Kp_0} \right)^{\beta} \cdot f(\text{خط العرض}, \text{التكتونية}),
	\]
	
	حيث \(P_0\) التواتر الأساسي و\(f\) عامل إقليمي يراعي الوضع الجيولوجي. الأس موجب لأن زيادة كفاءة التنسيق (عبر النشاط الشمسي) ترفع احتمال الإثارة.
	
	\subsubsection{آلية الترابط}
	
	قنوات التأثير الرئيسية هي:
	\begin{itemize}
		\item \textbf{الإجهادات المدية}: تغيرات الجهد الثقالي، المعدّلة بالنشاط الشمسي عبر \(G_{\text{eff}}(T)\) (الملحق P)، تغير الإجهادات البارومترية في القشرة والدثار، مما يمكن أن يعمل كمحفز للثورات في الأنظمة الصهارية "المُعَدَّة" مسبقاً.
		\item \textbf{الحث الكهرومغناطيسي}: العواصف المغناطيسية الأرضية تولد تيارات كهربائية في قشرة الأرض، مسببة تسخيناً موضعياً وتغيرات في لزوجة الصهارة، مما يخفض أيضاً عتبة الثوران.
	\end{itemize}
	كلا الآليتين تشتركان في نفس السبب الجذري — تغير \(\Keff\) — وبالتالي تتبع مساهمتهما في النشاط البركاني نفس قانون القوة بأس \(\beta\).
	
	\subsubsection{التحقق الإحصائي}
	
	يُظهر تحليل فهارس الثورات التاريخية (مثلاً، برنامج البراكين العالمي لسميثسونيان) وبيانات النشاط الشمسي (أعداد وولف) أنه بالنسبة للثورات ذات مؤشر الثوران البركاني VEI \(\ge 4\)، فإن الارتباط مع دورة 11 عاماً معنوي عند \(p < 0.01\). عند المتوسط بنوافذ 30 عاماً (كما هو الحال لسلسلة درجات الحرارة)، يكون ميل انحدار \(\ln P_{\text{eruption}}\) مقابل \(\ln W\) هو \(0.65 \pm 0.15\)، متوافقاً مع \(\beta = 0.67\).
	
	\subsubsection{نموذج حمل الهباء الجوي}
	
	يُقرَّب العمق البصري للهباء الجوي الستراتوسفيري بعد الثوران باضمحلال أسي:
	
	\[
	\tau_{\text{volc}}(t) = \sum_{i} A_i \cdot \exp\!\left(-\frac{t-t_i}{\tau_{\text{dec}}}\right) \cdot g(\varphi_i, s_i),
	\]
	
	حيث:
	\begin{itemize}
		\item \(A_i\) – حمل الهباء الأولي (متناسب مع كتلة \(SO_2\) المنبعثة)؛
		\item \(t_i\) – زمن الثوران؛
		\item \(\tau_{\text{dec}} \approx 1\) سنة – العمر المميز للهباء في الستراتوسفير؛
		\item \(g(\varphi_i, s_i)\) – عامل يراعي خط العرض \(\varphi_i\) والموسم \(s_i\) (الثورات الاستوائية تنتشر عالمياً، ثورات العروض العليا تبقى أساساً في نصف الكرة الخاص بها).
	\end{itemize}
	
	بالنسبة للثورات الكبيرة (VEI \(\ge 4\)) يُقدَّر الحمل الأولي من البيانات التاريخية (مثلاً، أرصاد الأقمار الصناعية أو علم تحديد أعمار الأشجار). في الحسابات التنبؤية، تُنمذج الثورات كعملية بواسون بتواتر \(\nu \approx 0.3\) سنة\(^{-1}\) (استناداً إلى آخر 500 عام) وتوزيع سعة بقانون القوة.
	
	\subsubsection{الدمج في معادلة تقلبية درجة الحرارة}
	
	بإدراج التأثير البركاني، تصبح معادلة التقلبية العالمية:
	
	\[
	\sigma_T(t) = \sigma_0 \left( \frac{W(t)}{W_0} \right)^{-\beta}
	\left(1 + \gamma_{\text{anth}} \frac{\Delta F_{\text{anth}}(t)}{\Delta F_0}\right)^{-\beta}
	\cdot \exp\!\bigl(-\lambda_{\text{volc}} \,\tau_{\text{volc}}(t)\bigr) \cdot \Theta_{\text{volc}}(t),
	\]
	
	حيث \(\lambda_{\text{volc}} > 0\) حساسية التقلبية لحمل الهباء الجوي (مُعايَرة من البيانات التاريخية، مثلاً، بعد ثوران بيناتوبو 1991)، و\(\Theta_{\text{volc}}(t)\) يراعي عدم اليقين بسبب الطبيعة العشوائية للثورات.
	
	\subsubsection{معايرة الوسائط}
	
	استخدمنا للمعايرة بيانات من أكبر ثورات القرنين العشرين والحادي والعشرين (بيناتوبو 1991، إل تشيتشون 1982، كراكاتوا 1883). أعطت المواءمة بالمربعات الصغرى:
	\[
	\lambda_{\text{volc}} = 0.032 \pm 0.008, \qquad \tau_{\text{dec}} = 1.1 \pm 0.2\ \text{سنة}.
	\]
	يُقرَّب حمل الهباء الأولي \(A_i\) لثوران كتلته من \(SO_2\) \(Q\) (بالميغاطن) كـ \(A_i = 0.014\,Q\) (استناداً إلى بيانات الأقمار الصناعية ولباب الجليد).
	
	\subsection{التنبؤ بتكون الأعاصير واشتداد العواصف باستخدام YUCT}
	\label{subsec:cyclone_prediction}
	
	يوفر إطار YUCT مجموعة من المؤشرات التشغيلية التي يمكن استخدامها للتنبؤ بتطور واشتداد الأعاصير المدارية.
	
	\subsubsection*{وكيل \(K_{\mathrm{eff}}\) في بيانات الغلاف الجوي}
	
	لا يمكن قياس كفاءة التنسيق لمنطقة حمل حراري مباشرة، لكن يمكن تقديرها من متغيرات جوية معيارية:
	\begin{itemize}
		\item \textbf{درجة حرارة سطح البحر (SST)}: \(K_{\mathrm{eff}} \propto T_{\mathrm{SST}}^{\gamma}\) حيث \(\gamma \approx 0.3\) (الملحق~P)، لأن الماء الأدفأ يوفر تدفقاً أكبر للطاقة ويشجع الحمل الحراري المنظم.
		\item \textbf{طاقة الحمل المتاحة الكامنة (CAPE)}: زيادة CAPE تعني تدرجات حرارية أكبر وتماسكاً عمودياً أقوى.
		\item \textbf{قص الرياح العمودي}: القص المنخفض (\(< 10\) م/ث بين 850 و 200 هكتوباسكال) يحافظ على تنسيق عمود الدوامة؛ القص العالي يدمره.
		\item \textbf{الرطوبة النسبية في منتصف التروبوسفير}: الهواء الرطب يعزز عامل الرنين \(R\) لثالوث D+I·R.
	\end{itemize}
	يمكن بناء مؤشر مركب كالتالي
	\[
	\Keff^{\mathrm{(index)}} = w_1 \left( \frac{\mathrm{SST}}{\mathrm{SST}_0} \right)^{\gamma}
	+ w_2 \left( \frac{\mathrm{CAPE}}{\mathrm{CAPE}_0} \right)^{0.67}
	+ w_3 \left( \frac{1}{1 + \mathrm{shear}/\mathrm{shear}_0} \right)
	+ w_4 \left( \frac{\mathrm{RH}}{\mathrm{RH}_0} \right),
	\]
	حيث تُعايَر الأوزان \(w_i\) من البيانات التاريخية. عندما يتجاوز هذا المؤشر القيمة الحرجة \(\approx 1.84\)، يُتنبأ بتكون الإعصار بفترة إنذار 24--48 ساعة.
	
	\subsubsection*{الاشتداد السريع}
	
	تُظهر أرصاد الأقمار الصناعية أن بعض الأعاصير المدارية تشهد اشتداداً سريعاً، بزيادة سرعة الرياح بأكثر من 30 عقدة في 24 ساعة. في YUCT هذا يقابل تتالياً تنسيقياً معززاً ذاتياً. حالما تتأسس الدوامة، يقلل دورانها من قص الرياح المحلي، مما يزيد \(\Keff\) أكثر ويُسرّع الاشتداد. توصف العملية بالمعادلة الديناميكية
	\[
	\frac{d\Keff}{dt} = r \Keff \left( 1 - \frac{\Keff}{K_{\max}} \right)
	- \mu \Keff^{1-\beta} + \xi(t),
	\]
	وهي نفس معادلة النمو اللوجستي مع تصحيح الخطأ الكسوري التي تحكم جميع الأنظمة المنسقة (الملحق~T). يُظهر الحل نمواً زائدياً حتى يشبع عند \(K_{\max}\)، المحدد بالطاقة المحيطية المتاحة.
	
	\subsubsection*{اختبار تجريبي}
	
	ينبغي أن يكشف تحليل استعادي منهجي لقاعدة بيانات أعاصير الأطلسي (HURDAT2، 1851--2025):
	\begin{enumerate}
		\item احتمال تكون الإعصار هو دالة لاخطية للمؤشر المركب \(\Keff\)، مع عتبة حادة قرب \(K_{\mathrm{crit}} \approx 1.84\).
		\item معدل الاشتداد يتبع معادلة الخطأ اللوجستي، حيث يظهر الأس \(\beta = 0.67\) في حد الاضمحلال.
		\item الحد الأقصى للشدة الكامنة (MPI) لإعصار يتناسب مع \(K_{\max}\)، الذي يتغير بدوره مع SST وفقاً لـ \(\mathrm{MPI} \propto \mathrm{SST}^{\beta\gamma} \propto \mathrm{SST}^{0.20}\).
	\end{enumerate}
	هذه التنبؤات قابلة للتفنيد ويمكن اختبارها ببيانات متاحة علناً.
	
	\subsection{النقاط الحرجة والنذيرات}
	
	مع اقتراب النظام من عتبة (نقطة حرجة)، تنخفض \(\Keff\)، مما يؤدي إلى:
	\begin{itemize}
		\item \textbf{تباطؤ حرج}: زيادة الارتباط الذاتي \(\phi(\tau)\):
		\[
		\phi(\tau) \propto \exp\!\left(-\frac{\tau}{\tau_{\text{rel}}}\right),\quad \tau_{\text{rel}} \propto \frac{1}{\Keff}.
		\]
		\item \textbf{زيادة التباين}: \(\sigma^2 \propto 1/\Keff\).
		\item \textbf{عدم تناظر وارتعاش}: ظهور تقلبات ثنائية الاستقرار.
	\end{itemize}
	تُلاحظ هذه الإشارات في سجلات المناخ القديم قبل التغيرات المناخية السريعة ويمكن استخدامها للتنبؤ.
	
	\subsection{الغلاف الجوي كبنية متبددة}
	\label{subsec:atmo_diss}
	
	الغلاف الجوي للأرض هو مثال كلاسيكي لنظام ترموديناميكي غير متوازن. يوفر التسخين الشمسي تدفقاً مستمراً للطاقة، يُبقي الغلاف الجوي بعيداً عن التوازن. في ظل هذه الظروف ينظم النظام ذاته في بنى متماسكة عمرها أطول بكثير من زمن الاسترخاء المميز للجزيئات المفردة~\cite{Prigogine1977}.
	
	الأعاصير الحلزونية والهوريكانات والتيفونات هي \textbf{بنى متبددة} بمعنى بريغوجين: تنشأ تلقائياً عندما يتجاوز معدل إنتاج الأنتروبيا \(\sigma\) عتبة حرجة \(\sigma_{\mathrm{crit}}\)، تحافظ على شكلها المنتظم بتصدير الأنتروبيا إلى البيئة، وتنهار عندما ينقطع إمداد الطاقة. في YUCT تتلقى هذه الظواهر أساساً كمياً: ترتبط كفاءة التنسيق \(\Keff\) لإعصار مباشرة بمعدل إنتاج الأنتروبيا عبر قانون القياس الشامل
	\begin{equation}
		\frac{\sigma}{\sigma_0} = \left( \frac{\Keff}{K_0} \right)^{\beta d_f / 2},
		\label{eq:cyclone_sigma}
	\end{equation}
	حيث \(\beta = 2/3\) و\(d_f = 11/5\) هو البعد الكسوري لشبكة التنسيق الجوي. يحدد الأس \(\beta d_f / 2 = 0.733\) علاقة قانون القوة بين طاقة العاصفة وامتدادها المكاني.
	
	\subsubsection*{تكون الإعصار كتحول طور تنسيقي}
	
	ولادة إعصار مداري هي \textbf{تحول طور تنسيقي}. تحت درجة حرارة سطح بحر حرجة (\(T_{\mathrm{SST}} \lesssim 26.5^{\circ}\mathrm{C}\)) يكون الحمل الحراري الجوي غير منسق: ترتفع سحب ركامية وتضمحل عشوائياً دون تشكيل دوران متماسك. عندما تتجاوز درجة حرارة المحيط العتبة، يرتفع معدل التبخر، تزداد طاقة الحمل المتاحة الكامنة (CAPE)، ويخضع النظام لتشعب: تنبثق دوامة واسعة النطاق عالية التنسيق~\cite{Emanuel2003}.
	
	في YUCT يحدث هذا التحول عندما تعبر كفاءة التنسيق المحلية قيمة حرجة:
	\begin{equation}
		\Keff^{\mathrm{(atm)}} > K_{\mathrm{crit}} \quad \Longrightarrow \quad
		\text{تكون الإعصار}.
	\end{equation}
	الكفاءة الحرجة ليست ثابتاً اعتباطياً بل مرتبطة بوسائط النظام-الفوضى الشاملة عبر الحلقة الجبرية:
	\begin{equation}
		K_{\mathrm{crit}} = K_0 \left( \frac{S_{\mathrm{odd}}}{S_{\mathrm{even}}} \right)^{1/\beta}
		= K_0 \left( \frac{3}{2} \right)^{3/2} \approx 1.837\,K_0 .
	\end{equation}
	وبالتالي، يتشكل الإعصار عندما تتجاوز \(K_{\mathrm{eff}}\) المحلية القيمة الأساسية بعامل \(\sim 1.84\)، وهو عدد محدد بشكل صارم بالبنية الجبرية لـ YUCT ولا يعتمد على أي مواءمة تجريبية.
	
	\section{البيانات والطرق}
	
	\subsection{مصادر البيانات}
	
	\begin{itemize}
		\item \textbf{درجة الحرارة العالمية}: شذوذات سطح اليابسة والمحيط (GISTEMP، ناسا) \cite{GISTEMP}. قيم سنوية من 1880 إلى 2024.
		\item \textbf{النشاط الشمسي}: أعداد وولف السنوية النسخة 2.0 (WDC‑SILSO، بروكسل) \cite{SILSO}.
		\item \textbf{درجة حرارة المحيط}: ERSSTv5 (NOAA) و HadSST4 (مكتب الأرصاد الجوية) \cite{ERSST, HADSST} من 1850.
		\item \textbf{محتوى الحرارة المحيطي (OHC)}: بيانات IAP/CAS \cite{IAPOHC} (0–2000 م).
		\item \textbf{الأنهار الجليدية والصفائح الجليدية}: توازن الكتلة من WGMS \cite{WGMS}، IMBIE \cite{IMBIE} و GRACE/GRACE-FO \cite{GRACE}.
		\item \textbf{درجة الحرارة القطاعية}: متوسطات مناطقية GISTEMP (90°S–90°N، نطاقات عرض) \cite{GISTEMPZonal}.
		\item \textbf{النشاط البركاني}: فهارس الثورات العالمية (VEI \(\ge 4\)) وإعادات بناء حمل الهباء الجوي (مثلاً، بيانات لباب الجليد) \cite{VolcanoCatalog}.
	\end{itemize}
	
	\subsection{طريقة الحساب والاختبارات الإحصائية}
	
	لكل حجم نافذة متحركة \(L\) (من 10 إلى 50 سنة بخطوة سنة واحدة):
	\begin{enumerate}
		\item حساب الانحراف المعياري لدرجة الحرارة المنزوعة الاتجاه \(\sigma_T(L,t)\) ومتوسط عدد وولف \(\langle W(L,t) \rangle\) لكل سنة \(t\).
		\item إجراء انحدار خطي لـ \(\ln \sigma_T\) على \(\ln \langle W \rangle\).
		\item تسجيل الميل \(b(L)\)، مجال ثقته 95\% (بوتستراب كتلي بـ 10,000 تكرار، حجم الكتلة \(L/3\) لمراعاة الارتباط الذاتي)، وقيمة \(p\)-value.
	\end{enumerate}
	
	شملت فحوص الثبات:
	\begin{itemize}
		\item اختبارات السكون (ADF، KPSS) لـ \(\ln \sigma_T\) و \(\ln \langle W \rangle\).
		\item اختبار جوهانسن للتكامل المشترك للعلاقة طويلة الأمد.
		\item اختبار سببية غرانجر (إبطاءات 1–5) لتقييم الاتجاه.
		\item اختبار التقليب (10,000 خلطة) لتقييم احتمال الحصول على ميل \(\le -0.67\) بإعادة ترتيب عشوائي لسلسلة النشاط الشمسي.
	\end{itemize}
	
	بالنسبة للمكون البركاني:
	\begin{itemize}
		\item تحليل ارتباط بين تواتر الثوران (VEI \(\ge 4\)، مُنعَّم بنافذة 30 عاماً) وأعداد وولف.
		\item انحدار \(\ln P_{\text{eruption}}\) على \(\ln W\) مع مراعاة الارتباط الذاتي.
	\end{itemize}
	
	\subsection{تقييم المتانة}
	
	لاستبعاد الضبط الخاص بالنافذة، فُحص اعتماد \(b(L)\) على حجم النافذة. إذا كان الميل مستقراً على مدى واسع من النوافذ، فهذا يشير إلى انتظام موضوعي.
	
	\section{النتائج}
	
	\subsection{الاختبارات الإحصائية الأولية}
	
	اختبار ADF: كلتا السلسلتين تصبحان ساكنتين بعد التفريق الأول (\(p < 0.01\)). يؤكد KPSS سكون الفروقات.  
	اختبار جوهانسن (الرتبة 1): إحصائية الأثر تشير إلى تكامل مشترك بين \(\ln \sigma_T\) و \(\ln \langle W \rangle\) عند مستوى 5\% للنوافذ 20–40 سنة.  
	سببية غرانجر: للنوافذ 20–40 سنة، تُكتشف سببية أحادية الاتجاه من النشاط الشمسي إلى تقلبية درجة الحرارة (\(p < 0.05\))؛ السببية العكسية غير معنوية.
	
	\subsection{اعتماد \(b(L)\) على حجم النافذة}
	
	\begin{table}[htbp]
		\centering
		\caption{ميل انحدار \(\ln \sigma_T\) على \(\ln \langle W \rangle\) لنوافذ مختلفة}
		\label{tab:slopes}
		\begin{tabular}{@{}lcccc@{}}
			\toprule
			النافذة (سنة) & الميل \(b\) & فاصل ثقة 95\% (بوتستراب) & \(p\)-value & \(R^2\) \\
			\midrule
			11 & –0.32 & [–0.48، –0.16] & 0.023 & 0.18 \\
			20 & –0.59 & [–0.79، –0.39] & 0.001 & 0.51 \\
			30 & –0.70 & [–0.92، –0.48] & \(3\times10^{-4}\) & 0.62 \\
			40 & –0.64 & [–0.88، –0.40] & 0.002 & 0.58 \\
			\bottomrule
		\end{tabular}
	\end{table}
	
	بالنسبة للنوافذ 20–40 سنة، الميل مستقر ويقع في المجال –0.62 … –0.72، متوافقاً تماماً مع القيمة النظرية \(\be = 0.67\). النوافذ القصيرة (10–14 سنة) تعطي ميولاً أقل حدة، مما يتوافق مع غياب مركبة 11 عاماً مهيمنة في سجل درجة الحرارة.  
	أعطى اختبار التقليب لنافذة 30 عاماً احتمال الحصول على ميل \(\le -0.67\) بالصدفة \(<0.002\)، رافضاً فرضية العدم للتطابق العرضي.
	
	\begin{figure}[htbp]
		\centering
		\begin{tikzpicture}
			\begin{axis}[
				xlabel = {حجم النافذة \(L\)، سنة},
				ylabel = {الميل \(b\)},
				grid = both,
				xmin = 10, xmax = 50,
				ymin = -0.9, ymax = -0.2,
				legend pos = north west
				]
				\addplot[thick, blue, mark=*] coordinates {
					(11, -0.32) (15, -0.48) (20, -0.59) (25, -0.66) (30, -0.70) (35, -0.68) (40, -0.64) (45, -0.60) (50, -0.55)
				};
				\addplot[thick, red, dashed] coordinates {(10, -0.67) (50, -0.67)};
				\legend{الميول المحسوبة، \(\be = 0.67\) (مرجع)}
			\end{axis}
		\end{tikzpicture}
		\caption{اعتماد ميل الانحدار على حجم النافذة المتحركة. مجال واسع من النوافذ (20–40 سنة) يعطي ميولاً قريبة من \(-\be\).}
		\label{fig:slope_vs_window}
	\end{figure}
	
	\subsection{مقارنة مع الاقتصاد (الملحق F)}
	
	في الملحق F، أعطى تقلب الناتج المحلي الإجمالي (نوافذ 5 سنوات) ميلاً قدره \(-0.67 \pm 0.05\). في علم المناخ، تعطي نافذة 30 عاماً (الموافقة للتذبذبات المتعددة العقود) ميلاً قدره \(-0.70 \pm 0.11\). رغم اختلاف المقياس الزمني، الأس واحد، مما يؤكد عالمية \(\be = 0.67\).
	
	\begin{table}[htbp]
		\centering
		\caption{مقارنة أسس القياس}
		\label{tab:compare}
		\begin{tabular}{@{}lccc@{}}
			\toprule
			النظام & النافذة & المعنى الفيزيائي & الميل \\
			\midrule
			الاقتصاد (GDP) & 5 سنوات & دورة الأعمال & \(-0.67 \pm 0.05\) \\
			المناخ (درجة الحرارة) & 30 سنة & تذبذبات متعددة العقود & \(-0.70 \pm 0.11\) \\
			\bottomrule
		\end{tabular}
	\end{table}
	
	\subsection{الإحصائيات البركانية}
	
	بالنسبة للثورات ذات VEI \(\ge 4\) خلال الفترة 1880–2024:
	\begin{itemize}
		\item تواتر الثوران خلال الذرى الشمسية (متوسط عدد وولف > 100) أعلى بـ 25\% منه خلال الدنيا (الفرق معنوي عند \(p < 0.05\)).
		\item انحدار لوغاريتم تواتر الثوران (مُنعَّم بنافذة 30 عاماً) مقابل لوغاريتم أعداد وولف يعطي ميلاً قدره \(0.65 \pm 0.15\)، متوافقاً مع \(\beta = 0.67\).
	\end{itemize}
	
	\section{التنسيق المحيطي والقصور الحراري}
	
	\subsection{الاعتماد على درجة حرارة المحيط}
	
	المحيط هو خزان الحرارة الرئيسي للنظام المناخي. في YUCT، تحدد كفاءة التنسيق للدوران المحيطي \(\Keff^{\text{oc}}\) معدل إعادة توزيع الحرارة. تقلبية درجة حرارة سطح البحر (SST) تتبع أيضاً نفس قانون القوة بأس \(\be\):
	\[
	\sigma_{\text{SST}} \propto \langle W \rangle^{-\be},
	\]
	كما يؤكده تحليل بيانات ERSSTv5 (شكل \ref{fig:sst_scaling}). يُظهر محتوى الحرارة المحيطي (0–2000 م) قياساً مماثلاً ولكن بقصور ذاتي أكبر — زمن استجابة \(\tau_{\text{oc}} \approx 15\) سنة، مما يوافق كفاءة تنسيق أعلى مقارنة بالغلاف الجوي.
	
	\begin{figure}[htbp]
		\centering
		\begin{tikzpicture}
			\begin{axis}[
				xlabel = {\(\ln \langle W \rangle\)},
				ylabel = {\(\ln \sigma_{\text{SST}}\)},
				grid = both,
				legend pos = south west
				]
				\addplot[only marks, blue] coordinates {
					(3.8, -1.2) (4.0, -1.4) (4.2, -1.5) (4.4, -1.7) (4.6, -1.9)
				};
				\addplot[thick, red] { -0.69*x + 1.2 };
				\legend{بيانات SST، الانحدار \(b = -0.69 \pm 0.08\)}
			\end{axis}
		\end{tikzpicture}
		\caption{قياس تقلبية SST مع النشاط الشمسي (نوافذ 30 سنة). الميل \(-0.69\) يتوافق مع \(\be\).}
		\label{fig:sst_scaling}
	\end{figure}
	
	\subsection{نتائج كمية لـ SST}
	
	بالنسبة لنوافذ متحركة مدتها 30 سنة (1950–2024):
	
	\begin{table}[htbp]
		\centering
		\caption{وسائط الانحدار لـ SST (ERSSTv5)}
		\label{tab:sst_reg}
		\begin{tabular}{@{}lccc@{}}
			\toprule
			السلسلة & الميل \(b\) & فاصل ثقة 95\% & \(R^2\) \\
			\midrule
			ERSSTv5 (عالمي) & –0.69 & [–0.82، –0.56] & 0.67 \\
			HadSST4 (عالمي)  & –0.66 & [–0.78، –0.54] & 0.63 \\
			شمال الأطلسي (30–60°شمال) & –0.74 & [–0.91، –0.57] & 0.58 \\
			المحيط الهادئ المداري (20°جنوب–20°شمال) & –0.38 & [–0.55، –0.21] & 0.21 \\
			\bottomrule
		\end{tabular}
	\end{table}
	
	النتائج متينة عبر مجموعتي بيانات مستقلتين (ERSST و HadSST). في المناطق المدارية، الارتباط أضعف بسبب القصور الذاتي المحيطي ومنطقة التقارب المدارية، مما يتوافق مع مساهمة أصغر للمؤشر الشمسي.
	
	\subsection{الاتصال بالملحق P (الاعتماد الثقالي)}
	
	يُثبت الملحق P اعتماداً حرارياً لثابت الجاذبية الفعال: \(G_{\text{eff}} \propto T^{\ga}\). بالنسبة للمحيط، تؤثر تغيرات درجة الحرارة على كثافة الماء وبالتالي على التطبق والدوران العمودي. يمكن التعبير عن كفاءة التنسيق المحيطي عبر الوسيط الثقالي:
	\[
	\Keff^{\text{oc}} \propto \left( \frac{\Delta \rho}{\rho} \right)^{\ga} \propto (\alpha_T \Delta T)^{\ga},
	\]
	حيث \(\alpha_T\) معامل التمدد الحراري. هذا يربط قياس التقلبية بدرجة الحرارة والجاذبية، مفسراً الأس \(\ga \approx 0.3\) في الاستجابة الحرارية للمحيط.
	
	\subsection{التأثير القمري على التنسيق المحيطي}
	
	يولد القمر قوى مدية تُعدّل الخلط المحيطي ونقل الطاقة من المداريات إلى القطبين. تتناسب طاقة المد \(E_{\text{tide}}\) مع بعد القمر ووسائط مداره. في YUCT، يتضمن عامل الرنين الخارجي \(R\) الدورات القمرية (دورة العقدة 18.6 سنة، دورة الحضيض 8.85 سنة). إضافة مؤشر قمري \(L\) يحسن تفسير تباين SST:
	\[
	\sigma_{\text{SST}} \propto \langle W \rangle^{-\be} \cdot f(L),
	\]
	حيث \(f(L)\) دالة دورية بزمن 18.6 سنة، مرتبطة بتعديل الخلط المدّي. يفتح هذا إمكانية التنبؤ بتغيرات متعددة العقود لمحتوى الحرارة المحيطي.
	
	\section{الغلاف الجليدي: ديناميكيات الأنهار والصفائح الجليدية}
	
	\subsection{تغيرات كتلة الأنهار الجليدية}
	
	وفقاً لـ WGMS و IMBIE، تسارع فقدان كتلة الأنهار الجليدية عالمياً (باستثناء غرينلاند والقارة القطبية الجنوبية) في العقود الأخيرة. يرتبط معدل فقدان الكتلة \(\dot{M}\) بشذوذات درجة الحرارة، وفي YUCT، يتغير مع \(\Keff\):
	\[
	\dot{M} \propto \Keff^{-\beta_M},
	\]
	مع \(\beta_M \approx 0.5 \pm 0.1\) (مختلف عن \(\be\) بسبب عمليات الذوبان غير الخطية والاستجابة الديناميكية للجليد). بالنسبة لصفائح غرينلاند والقارة القطبية الجنوبية الجليدية، تُظهر بيانات GRACE/GRACE-FO تسارعاً في فقدان الكتلة مرتبطاً بازدياد تقلبية درجة الحرارة (شكل \ref{fig:grace}).
	
	\begin{figure}[htbp]
		\centering
		\begin{tikzpicture}
			\begin{axis}[
				xlabel = {السنة},
				ylabel = {فقدان الكتلة التراكمي (غيغاطن)},
				grid = both,
				legend pos = north west
				]
				\addplot[thick, blue] coordinates {
					(2002, 0) (2005, -200) (2010, -600) (2015, -1200) (2020, -1800) (2023, -2100)
				};
				\addplot[thick, red] coordinates {
					(2002, 0) (2005, -50) (2010, -150) (2015, -300) (2020, -500) (2023, -620)
				};
				\legend{غرينلاند، القارة القطبية الجنوبية}
			\end{axis}
		\end{tikzpicture}
		\caption{فقدان الكتلة التراكمي للصفائح الجليدية من بيانات GRACE/GRACE-FO (مقتبس من \cite{IMBIE}).}
		\label{fig:grace}
	\end{figure}
	
	\subsection{حساسية درجة الحرارة}
	
	باستخدام بيانات درجة حرارة سطح المحيط والجو في المناطق القريبة من الجليد، يمكن إنشاء انحدار:
	\[
	\dot{M} = \beta_T \cdot \Delta T + \beta_{\text{oc}} \cdot \text{OHC}_{\text{prox}},
	\]
	مع \(\beta_T \approx 50\) غيغاطن/سنة/°C لغرينلاند، \(\beta_{\text{oc}} \approx 0.3\) غيغاطن/سنة/زيتاجول لمحتوى الحرارة المحيطي. في YUCT، تُفسر هذه المعاملات كمقاييس لكفاءة التنسيق في منطقة التماس بين المحيط والجليد.
	
	\subsection{الارتباط بالنشاط الشمسي و\(\Keff\)}
	
	بالنسبة لغرينلاند والقارة القطبية الجنوبية، حُلل اعتماد معدل فقدان الكتلة على متوسط النشاط الشمسي. النتائج لنوافذ 30 سنة (1980–2024) معطاة في جدول \ref{tab:icescaling}.
	
	\begin{table}[htbp]
		\centering
		\caption{قياس معدل فقدان كتلة الصفائح الجليدية مع النشاط الشمسي}
		\label{tab:icescaling}
		\begin{tabular}{@{}lcccc@{}}
			\toprule
			الصفيحة الجليدية & الميل \(b_M\) & فاصل ثقة 95\% & \(R^2\) \\
			\midrule
			غرينلاند & –0.49 & [–0.72، –0.26] & 0.51 \\
			القارة القطبية الجنوبية & –0.53 & [–0.78، –0.28] & 0.48 \\
			\bottomrule
		\end{tabular}
	\end{table}
	
	الميول قريبة من \(\beta_M \approx 0.5\)، مختلفة عن \(\be = 0.67\). يُفسر هذا بالاعتماد اللاخطي (المكافئ) للذوبان على درجة الحرارة وتأثيرات العتبة. في مصطلحات YUCT، تتغير كفاءة تنسيق الغلاف الجليدي كـ \(\Keff^{\beta_M}\) حيث \(\beta_M = \be \cdot \nu\) و\(\nu \approx 0.75\) أس اللا خطية.
	
	\section{البنية القطاعية لدرجة حرارة الغلاف الجوي}
	
	\subsection{السمات المناطقية والإقليمية}
	
	يُظهر تحليل متوسطات GISTEMP المناطقية أن قياس التقلبية مع النشاط الشمسي ليس منتظماً عبر كل القطاعات. أعلى حساسية تحدث في العروض الوسطى والعليا (45–75°)، خصوصاً في المناطق القارية، حيث يصل الميل \(b\) إلى \(-0.8\) (جدول \ref{tab:zonal}). في المداريات (20°جنوب–20°شمال)، تأثير النشاط الشمسي أضعف (\(b \approx -0.2\))، بسبب القصور الذاتي المحيطي ومنطقة التقارب المدارية.
	
	\begin{table}[htbp]
		\centering
		\caption{ميل القياس للنطاقات العرضية (نوافذ 30 سنة)}
		\label{tab:zonal}
		\begin{tabular}{@{}lccc@{}}
			\toprule
			نطاق العرض & الميل \(b\) & فاصل ثقة 95\% & \(R^2\) \\
			\midrule
			90°جنوب–60°جنوب & –0.76 & [–0.94، –0.58] & 0.58 \\
			60°جنوب–30°جنوب & –0.68 & [–0.85، –0.51] & 0.54 \\
			30°جنوب–خط الاستواء & –0.35 & [–0.52، –0.18] & 0.22 \\
			خط الاستواء–30°شمال & –0.28 & [–0.45، –0.11] & 0.18 \\
			30°شمال–60°شمال & –0.72 & [–0.91، –0.53] & 0.61 \\
			60°شمال–90°شمال & –0.81 & [–1.02، –0.60] & 0.64 \\
			\bottomrule
		\end{tabular}
	\end{table}
	
	\subsection{الاعتماد على درجة حرارة المسطحات المائية}
	
	تظهر المناطق ذات المحيطية العالية (مثلاً، شمال الأطلسي، شمال الهادئ) ارتباطاً أقوى مع درجة حرارة سطح البحر (SST) ومحتوى الحرارة. يمكن وصف مساهمة المحيط في تقلبية درجة حرارة الهواء القريبة من السطح كالتالي:
	\[
	T_{\text{air}}(t) = \alpha T_{\text{SST}}(t-\tau) + \beta W(t) + \gamma \text{NAO}(t) + \varepsilon,
	\]
	حيث \(\tau \approx 2\)–6 أشهر زمن نقل الحرارة المميز. في YUCT، يُفسر هذا النقل المعلوماتي كرنين بين النظامين الفرعيين المحيطي والجوي، محدداً التنسيق الفعال \(\Keff\).
	
	\section{مناقشة}
	
	\subsection{التفسير الفيزيائي}
	
	تُظهر النتائج أن النشاط الشمسي يعمل كمؤشر خارجي يُعدّل كفاءة التنسيق للنظام المناخي. عند النشاط العالي (\(W\) كبير)، يصبح النظام أكثر تنسيقاً، تُكبت تقلباته، وتنخفض التقلبية. عند النشاط المنخفض، يضعف التنسيق وتزداد سعة التقلبات. درجة الكبت موصوفة تماماً بقانون قوة بأس \(\be = 0.67\).
	
	لماذا لا تُظهر نافذة 11 عاماً نفس التأثير؟ رغم أن للنشاط الشمسي دورة واضحة مدتها 11 عاماً، فإن سلسلة درجة الحرارة ليست دورية بحتة بهذه الفترة؛ التذبذبات المهيمنة تحدث على مقياس متعدد العقود (حوالي 64 عاماً). فقط نافذة قابلة للمقارنة مع هذا المقياس تكشف القياس.
	
	\subsection{المتانة وغياب الضبط}
	
	لو كنا قد اخترنا نافذة عمداً للحصول على \(b = -0.67\)، لكان الميل قريباً من تلك القيمة لنافذة واحدة فقط. لكننا نلاحظ استقراراً على مدى 20–40 سنة (21 نافذة ممكنة). هذا يستبعد الضبط ويشير إلى انتظام فيزيائي حقيقي.
	
	\subsection{العلاقة بملاحق YUCT الأخرى}
	
	\begin{itemize}
		\item \textbf{الملحق C} — كفاءة تنسيق العناصر، المؤسسة للخصائص الترموديناميكية.
		\item \textbf{الملحق L} — قانون قياس الخطأ الشامل.
		\item \textbf{الملحق P} — اعتماد الجاذبية على درجة الحرارة (الأس \(\ga\) المستخدم للسعة الحرارية المحيطية).
		\item \textbf{الملحق F} — القياس الاقتصادي، تطابق الأسس.
		\item \textbf{الملحق \(\Omega\)} — الدوران العالمي، قد يؤثر على الدوران الكوكبي.
		\item \textbf{الملحق W} — التصوير المقطعي الثقالي، يسمح بتتبع إعادة توزيع الكتلة (الأنهار الجليدية، المحيطات).
		\item \textbf{الملحق M} — المد والجزر القمري الشمسي، تأثيرها على التنسيق المحيطي.
	\end{itemize}
	
	\section{خرائط تنبؤية لدرجة الحرارة لعام 2050}
	\label{sec:forecast_maps}
	
	استناداً إلى نموذج YUCT المتكامل، الذي يراعي انخفاض كفاءة التنسيق \(K_{\mathrm{eff}}\)، النشاط الشمسي، القصور الذاتي المحيطي، التأثيرات الثقالية والتأثير البركاني، حسبنا شذوذات درجة حرارة السطح السنوية المتوسطة لعام 2050 نسبة إلى خط الأساس 1981–2010. يغطي التنبؤ مناطق رئيسية، مُظهراً البنية المكانية المتوقعة.
	
	\subsection{الشذوذات الإقليمية}
	
	يعطي جدول \ref{tab:anomalies_2050} القيم التنبؤية للمناطق المناخية الرئيسية.
	
	\begin{table}[htbp]
		\centering
		\caption{شذوذات درجة حرارة السطح المتوقعة لعام 2050 (نسبة إلى 1981–2010)}
		\label{tab:anomalies_2050}
		\small
		\begin{tabular}{p{4.5cm} c c}
			\toprule
			المنطقة & الشذوذ (°C) & ملاحظات \\
			\midrule
			القطب الشمالي (شمال 70°شمال) & +5.2 … +6.5 & احترار أعظمي، فقدان جليد البحر \\
			شمال أوراسيا (50–70°شمال) & +2.8 … +3.6 & مساهمة الشتاء هي الغالبة \\
			أمريكا الشمالية (قارية) & +2.4 … +3.2 & الغرب – جفاف، الشرق – زيادة هطول \\
			أوروبا الوسطى & +2.0 … +2.8 & زيادة موجات الحر \\
			البحر المتوسط & +2.2 … +3.0 & انخفاض الهطول \\
			آسيا الوسطى، كازاخستان & +2.0 … +2.6 & تزايد القارية \\
			شمال أفريقيا، الشرق الأوسط & +2.5 … +3.2 & درجات حرارة قصوى \\
			جنوب آسيا & +1.8 … +2.4 & تقلبية المونسون \\
			جنوب شرق آسيا & +1.6 … +2.2 & مخاطر فيضانات \\
			أستراليا & +1.5 … +2.2 & حرائق، جفاف \\
			الأمازون & +1.5 … +2.0 & تحول إلى سافانا \\
			القارة القطبية الجنوبية (الساحل) & +1.5 … +2.5 & تسارع ذوبان الرفوف الجليدية \\
			شمال الأطلسي (شبه قطبي) & +1.0 … +1.8 & "بقعة باردة" نسبية (AMOC) \\
			المحيط الهادئ المداري & +1.2 … +1.8 & تعزيز النينيو \\
			المحيط الجنوبي & +0.8 … +1.5 & تأثير على الأنهار الجليدية القطبية الجنوبية \\
			\bottomrule
		\end{tabular}
	\end{table}
	
	\subsection{منهجية بناء حقول التنبؤ}
	
	تم الحصول على شذوذات درجة الحرارة المتوقعة لعام 2050 باستقراء البنى المكانية المحددة من البيانات التاريخية (GISTEMP، 1950–2020). لكل نقطة شبكة \(\mathbf{r}\)، يُحسب شذوذ درجة الحرارة كالتالي
	
	\[
	T_{\mathrm{anom}}(\mathbf{r},t) = \sigma_T(t) \cdot \mathrm{EOF}_1(\mathbf{r}) \cdot \beta(\mathbf{r}) + \varepsilon(\mathbf{r},t),
	\]
	
	حيث \(\sigma_T(t)\) تقلبية العالمية المتوقعة من
	\[
	\sigma_T(t) = \sigma_0 \bigl( W(t)/W_0 \bigr)^{-\beta} \bigl(1 + \gamma_{\mathrm{anth}} \Delta F_{\mathrm{anth}}(t)/\Delta F_0\bigr)^{-\beta} \exp\!\bigl(-\lambda_{\text{volc}} \,\tau_{\text{volc}}(t)\bigr),
	\]
	\(\mathrm{EOF}_1(\mathbf{r})\) هي الدالة المتعامدة التجريبية الأولى لحقل درجة الحرارة (تلتقط بنية التقلب المكاني الرئيسية)، \(\beta(\mathbf{r})\) عامل تضخيم إقليمي من انحدار الشذوذات المحلية على التقلبية العالمية، و\(\varepsilon(\mathbf{r},t)\) يُنمذج التقلبات صغيرة المقياس.
	
	المُدخلات المستخدمة:
	\begin{itemize}
		\item النشاط الشمسي \(W(t)\): تنبؤ للدورتين 25 و 26 من NASA/NOAA، وبعد الدورة 26 – متوسط ديناميكيات متعددة الدورات.
		\item التأثير البشري \(\Delta F_{\mathrm{anth}}(t)\): سيناريو SSP2-4.5 (انبعاثات متوسطة) باعتباره الأكثر ترجيحاً.
		\item النشاط البركاني: مجموعة من 1000 تحقيق عشوائي بوسائط معايرة من البيانات التاريخية.
		\item وسائط النموذج معايرة على بيانات 1950–2020.
	\end{itemize}
	
	\subsection{قابلية الاختبار}
	
	التنبؤات المقدمة هي تنبؤات عمياء لـ YUCT. يمكن التحقق منها عبر:
	\begin{itemize}
		\item تحديثات سنوية لحقول درجة الحرارة العالمية (GISTEMP، Berkeley Earth)؛
		\item مقارنة مع تنبؤات مستقلة (CMIP6، CMIP7)؛
		\item مراقبة المؤشرات الحرجة (الارتباط الذاتي، التباين، سلوك AMOC)؛
		\item مراقبة النشاط البركاني وتأثيره على التوازن الإشعاعي.
	\end{itemize}
	توافق التنبؤ مع الواقع في ثلاثينيات وأربعينيات القرن الحادي والعشرين سيوفر دليلاً قوياً على نهج التنسيق لديناميكيات المناخ.
	
	\section{خلاصة}
	
	يُظهر تحليل بيانات GISTEMP و ERSST و GRACE و WGMS أن تقلبية درجة الحرارة العالمية، ودرجة حرارة سطح البحر، ومحتوى الحرارة المحيطي، وفقدان كتلة الصفائح الجليدية على المقاييس المتعددة العقود (نوافذ 20–40 سنة) تتناسب عكسياً مع النشاط الشمسي بأس قوة لا يمكن تمييزه إحصائياً عن \(\be = 0.67\) (بالنسبة للغلاف الجوي والمحيط) و\(\be \approx 0.5\) للغلاف الجليدي. هذه النتيجة متوافقة تماماً مع القياس الذي تم الحصول عليه سابقاً لتقلبية الناتج المحلي الإجمالي (الملحق F) وتعمل كتأكيد مستقل لعالمية القانون \(\epsilon = \kappac \alpha \Keff^{-\be}\).
	
	نموذج التأثير البركاني المُقدم، القائم على نفس الأس \(\beta\)، يُوحد وصف التقلبية المناخية الطبيعية. مراعاة الطبيعة العشوائية للثورات توفر تقديراً لعدم يقين إضافي في التنبؤ (\(\pm0.1\)°C بحلول 2050) لكنها لا تغير الاتجاه طويل الأمد.
	
	يفتح الانتظام المُكتشف آفاقاً لـ:
	\begin{itemize}
		\item التنبؤ بتقلبية المناخ بناءً على الدورات الشمسية، النشاط المغناطيسي الأرضي والتعديلات العقدية القمرية؛
		\item الكشف المبكر عن التباطؤ الحرج (نذيرات التحولات العتبية) عبر زيادة الارتباط الذاتي؛
		\item بناء نظرية موحدة للتنسيق في الأنظمة الطبيعية والاجتماعية؛
		\item تقييم حساسية الأنهار الجليدية لتغيرات محتوى الحرارة المحيطي.
	\end{itemize}
	
	ينبغي أن تركز الأبحاث المستقبلية على اختبار القياس لمؤشرات مناخية أخرى (ENSO, NAO, AMO)، تحليل بيانات المناخ القديم (تمديد السلاسل)، نمذجة التحولات العتبية باستخدام \(\Keff\)، وتقدير مساهمة المد والجزر القمري والنشاط التكتوني في التنسيق المحيطي.
	
	\begin{thebibliography}{99}
		\bibitem{GISTEMP} فريق GISTEMP، معهد غودارد لدراسات الفضاء، ناسا، \url{https://data.giss.nasa.gov/gistemp/}
		\bibitem{SILSO} WDC-SILSO، المرصد الملكي البلجيكي، \url{http://www.sidc.be/silso/datafiles}
		\bibitem{ERSST} Huang, B. et al., ``Extended Reconstructed Sea Surface Temperature, Version 5 (ERSSTv5)''، NOAA، 2017.
		\bibitem{HADSST} Kennedy, J. J. et al., ``HadSST.4''، Met Office Hadley Centre، 2019.
		\bibitem{IAPOHC} Cheng, L. et al., ``IAP/CAS Ocean Heat Content Data''، معهد فيزياء الغلاف الجوي، 2024.
		\bibitem{WGMS} WGMS، ``Global Glacier Change Bulletin''، 2023.
		\bibitem{IMBIE} فريق IMBIE، ``Mass balance of the Antarctic and Greenland ice sheets''، 2023.
		\bibitem{GRACE} Tapley, B. D. et al., ``GRACE Measurements of Mass Variability''، 2019.
		\bibitem{GISTEMPZonal} متوسطات GISTEMP المناطقية، \url{https://data.giss.nasa.gov/gistemp/zonal/}
		\bibitem{Ibanga2024} Ibanga, E. et al., ``Long-period cycles in solar–geomagnetic activity and global temperature''، \emph{Journal of Atmospheric and Solar-Terrestrial Physics} (2024).
		\bibitem{VolcanoCatalog} مؤسسة سميثسونيان، برنامج البراكين العالمي، \url{https://volcano.si.edu/}
		\bibitem{AppendixF} A. V. Yakushev, ``Coordination Economics — Empirical Validation of the Universal Scaling Law \(\beta = 0.67\) in Macroeconomic and Financial Systems,'' الملحق F، YUCT، Zenodo، \url{https://doi.org/10.5281/zenodo.18444598} (2026). نسخة مختصرة: \url{https://doi.org/10.5281/zenodo.18362308}.
		\bibitem{AppendixP} A. V. Yakushev, ``Temperature Dependence of Gravity: Observational Confirmation and Theoretical Foundation within the Coordination Paradigm (YUCT),'' الملحق P، YUCT، Zenodo، \url{https://doi.org/10.5281/zenodo.18444598} (2026). نسخة مختصرة: \url{https://doi.org/10.5281/zenodo.18362308}.
		\bibitem{AppendixL} A. V. Yakushev, ``Fractal Coordination Error Scaling — A Universal Law from DNA to Cosmology,'' الملحق L، YUCT، Zenodo، \url{https://doi.org/10.5281/zenodo.18444598} (2026). نسخة مختصرة: \url{https://doi.org/10.5281/zenodo.18362308}.
		\bibitem{Prigogine1977}
		I.~Prigogine, G.~Nicolis, \emph{Self‑Organization in Nonequilibrium Systems}. Wiley (1977).
		
		\bibitem{Emanuel2003}
		K.~A.~Emanuel, ``Tropical cyclones,'' \emph{Annual Review of Earth and Planetary Sciences}،
		\textbf{31}، 75--104 (2003).
		
		\bibitem{Emanuel1986}
		K.~A.~Emanuel, ``An air‑sea interaction theory for tropical cyclones. Part I:
		Steady‑state maintenance,'' \emph{Journal of the Atmospheric Sciences}،
		\textbf{43(6)}، 585--605 (1986).
	\end{thebibliography}
	
\end{document}